\documentclass[11pt,a4paper]{article}
% =====================================================================
%  Gemeinsame Vorlage der Arbeitsblätter "Graphen" – Informatik 11 (NTG)
%  (gleiches Layout wie informatik/10/java/arbeitsblaetter/ab-vorlage.tex)
%  Einbinden in jedem Blatt direkt nach \documentclass: % =====================================================================
%  Gemeinsame Vorlage der Arbeitsblätter "Graphen" – Informatik 11 (NTG)
%  (gleiches Layout wie informatik/10/java/arbeitsblaetter/ab-vorlage.tex)
%  Einbinden in jedem Blatt direkt nach \documentclass: % =====================================================================
%  Gemeinsame Vorlage der Arbeitsblätter "Graphen" – Informatik 11 (NTG)
%  (gleiches Layout wie informatik/10/java/arbeitsblaetter/ab-vorlage.tex)
%  Einbinden in jedem Blatt direkt nach \documentclass: \input{ab-vorlage}
%  Übersetzen mit XeLaTeX, LuaLaTeX, pdfLaTeX oder Tectonic.
%
%  Blatt:   xelatex ab-3-1-grundbegriffe.tex
%  Lösung:  xelatex ab-3-1-grundbegriffe-lsg.tex
%           (Datei enthält nur: \def\mitloesung{}\input{ab-3-1-grundbegriffe})
%           (füllt alle Lücken rot aus; siehe \lsg, \lsglinien, \lsgoder)
% =====================================================================
\usepackage[a4paper,left=18mm,right=18mm,top=26mm,bottom=17mm,headheight=15mm,headsep=4mm,footskip=8mm]{geometry}
\usepackage{iftex}
\ifPDFTeX
  \usepackage[T1]{fontenc}
  \usepackage[utf8]{inputenc}
  \usepackage{helvet}
  \usepackage[scaled=0.86]{beramono}
  \renewcommand{\familydefault}{\sfdefault}
\else
  \usepackage{fontspec}
  \setmainfont{texgyreheros}[Extension=.otf,UprightFont=*-regular,BoldFont=*-bold,ItalicFont=*-italic,BoldItalicFont=*-bolditalic]
  \setmonofont{DejaVuSansMono}[Extension=.ttf,UprightFont=*,BoldFont=*-Bold,ItalicFont=*-Oblique,BoldItalicFont=*-BoldOblique,Scale=0.86]
\fi
\usepackage[ngerman,shorthands=off]{babel}
\usepackage{xcolor,graphicx,tabularx,array,enumitem,fancyhdr,tikz,listings,multicol,calc,etoolbox}
\usepackage[most]{tcolorbox}
\usetikzlibrary{arrows.meta,positioning,calc,decorations.pathreplacing,shapes.geometric}
\usepackage[hidelinks]{hyperref}
\setlength{\parindent}{0pt}
\setlength{\parskip}{3pt}
\setlist{nosep,leftmargin=*}

% ---------- Lösungsschalter ----------
\newif\ifloesung
\ifdefined\mitloesung\loesungtrue\fi
\newcommand{\lsgfarbe}{\color{red!75!black}}

% ---------- Kopf- und Fußzeile (einheitliche Form) ----------
\newcommand{\lehrkraft}{Hau}
\newcommand{\themenbereich}{3. Graphen}
\newcommand{\abnummer}{}
\pagestyle{fancy}\fancyhf{}
\renewcommand{\headrulewidth}{0.4pt}
\fancyhead[L]{\small Klasse 11\\Datum:}
\fancyhead[C]{\small Themenbereich:\\\themenbereich}
\fancyhead[R]{\small \lehrkraft\ifloesung\\{\lsgfarbe\bfseries Lösung}\fi}
\fancyfoot[C]{\scriptsize edu-mrh.de\,·\,Informatik 11\,·\,Blatt \abnummer\ifloesung\ (Lösung)\fi\,·\,Seite \thepage}
\newcommand{\abtitel}[2]{\renewcommand{\abnummer}{#1}\begin{center}\Large\bfseries #1\quad\underline{#2}\end{center}\vspace{-1mm}}

% ---------- Boxen ----------
\newenvironment{aufgabe}[2][]{\begin{tcolorbox}[enhanced,breakable,colback=white,colframe=black,boxrule=0.7pt,arc=3mm,left=2mm,right=2mm,top=1.5mm,bottom=1.5mm,before skip=2.5mm,after skip=2.5mm]\textbf{Aufgabe #2}\ifstrempty{#1}{}{ \textit{#1}}\par\vspace{0.6mm}}{\end{tcolorbox}}
\newtcolorbox{merke}[1][Merke]{enhanced,breakable,colback=green!5,colframe=green!40!black,boxrule=0.8pt,arc=2mm,left=2mm,right=2mm,top=1.5mm,bottom=1.5mm,before skip=2.5mm,after skip=2.5mm,title={#1},fonttitle=\bfseries,coltitle=white}
\newtcolorbox{hinweis}{enhanced,breakable,colback=yellow!12,colframe=orange!70!black,boxrule=0.6pt,arc=2mm,left=2mm,right=2mm,top=1mm,bottom=1mm,before skip=2mm,after skip=2mm}
\newtcolorbox{ausblick}[1][Ausblick]{enhanced,breakable,colback=teal!5,colframe=teal!60!black,boxrule=0.8pt,arc=2mm,left=2mm,right=2mm,top=1.5mm,bottom=1.5mm,before skip=2.5mm,after skip=2.5mm,title={#1},fonttitle=\bfseries,coltitle=white}
\newcommand{\web}[1]{{\small\color{gray}\textit{Website: #1}}}

% ---------- Lücken, Linien, Karos ----------
\newcommand{\luecke}[1][3.5cm]{\rule[-0.5ex]{#1}{0.4pt}}
\newcommand{\linien}[1]{\par\vspace{1.5mm}\foreach \i in {1,...,#1}{\noindent\rule[0pt]{\linewidth}{0.3pt}\par\vspace{5.5mm}}\vspace{-3mm}}
\newcommand{\kariert}[1]{\par\vspace{1.5mm}\noindent\begin{tikzpicture}\draw[step=5mm,gray!55,thin] (0,0) grid (\linewidth-1pt,#1*5mm);\end{tikzpicture}\par}
\newcommand{\karobox}[2]{\begin{tikzpicture}\draw[step=5mm,gray!55,thin] (0,0) grid (#1,#2);\end{tikzpicture}}
\newcommand{\code}[1]{\texttt{#1}}

% Lücke mit Lösung: \lsg[Mindestbreite]{Lösungstext}
\newlength{\lsgbreite}
\newcommand{\lsg}[2][3.5cm]{\ifloesung\settowidth{\lsgbreite}{\,#2\,}\ifdim\lsgbreite<#1\setlength{\lsgbreite}{#1}\fi\underline{\makebox[\lsgbreite][l]{\lsgfarbe\,#2}}\else\luecke[#1]\fi}
% Schreiblinien mit Lösung: \lsglinien{Anzahl}{Lösungstext} (gleiche Höhe wie \linien)
\newcommand{\lsglinien}[2]{\ifloesung\par\vspace{1.5mm}\noindent\begin{minipage}[t][#1\dimexpr 6.9mm\relax][t]{\linewidth}\lsgfarbe #2\end{minipage}\par\vspace{-1.5mm}\else\linien{#1}\fi}
% Beliebiger Inhalt: \lsgoder{nur in der Lösung}{nur im Blatt}
\newcommand{\lsgoder}[2]{\ifloesung\leavevmode#1\else#2\fi}

% ---------- Ergänzung (aus informatik/10/java/arbeitsblaetter/ab-vorlage.tex) ----------
% Bestehende Makros oben bleiben unverändert, damit ältere Blätter gleich aussehen.
%  \restzeile / \lsgrest{Text}   Lücke bis zum Zeilenende
%  \lsgtext{Text}                nur in der Lösung sichtbar (z. B. Tabellenzellen)
%  \lueckentext                  größerer Zeilenabstand, Flattersatz für Absätze mit Lücken
%  \schreibzeile                 Strut für Tabellenzellen zum Ausfüllen (8 mm hoch)
%  \alt{9}, \neu{7}              überschriebener alter Wert / neuer Wert (Dijkstra-Tabellen)
\usepackage{colortbl,xspace}
\usepackage[normalem]{ulem}
\newcommand{\lsgrest}[1]{\leavevmode\unskip\hspace{1.5mm}\ifloesung\rlap{\raisebox{0.8pt}{\lsgfarbe\,#1}}\fi\leaders\hbox{\rule[-0.6ex]{0.5pt}{0.4pt}}\hfill\null\par}
\newcommand{\restzeile}{\lsgrest{}}
\newcommand{\lsgtext}[1]{\ifloesung\leavevmode{\lsgfarbe#1}\fi}
\newcommand{\lueckentext}{\linespread{1.5}\selectfont\raggedright}
\newcommand{\schreibzeile}{\rule[-2.8mm]{0pt}{8mm}}
\newcommand{\alt}[1]{\leavevmode{\color{blue!60!black}\sout{#1}}}
\newcommand{\neu}[1]{\leavevmode{\color{green!50!black}\bfseries #1}}
\colorlet{djaktuell}{orange!40}
\colorlet{djfertig}{blue!22}

% ---------- Graphen zeichnen ----------
\tikzset{
  knoten/.style={draw,rounded corners=3pt,fill=orange!12,minimum height=6mm,inner sep=3pt,font=\small},
  kknoten/.style={draw,circle,fill=orange!12,minimum size=7mm,inner sep=1pt,font=\small},
  kante/.style={thick},
  gkante/.style={thick,-{Stealth[length=2.6mm]}},
  dkante/.style={thick,{Stealth[length=2.6mm]}-{Stealth[length=2.6mm]}},
  gewicht/.style={font=\scriptsize,fill=white,inner sep=1pt,text=teal!60!black},
  lsgkante/.style={thick,red!75!black},
  lsggkante/.style={thick,red!75!black,-{Stealth[length=2.6mm]}},
  lsgknoten/.style={knoten,draw=red!75!black,text=red!75!black,fill=red!4},
}
% Adjazenzmatrix zum Ausfüllen: \adjmx[Lösungseinträge]{Anzahl}{Knotenliste}{Zellbreite}
% Lösungseinträge als Zeile/Spalte/Text, z. B. \adjmx[1/2/X,2/1/X]{3}{A,B,C}{7mm}
\newcommand{\adjmx}[4][]{\begin{tikzpicture}[x=#4,y=#4]
  \foreach \n [count=\i] in {#3} {\node[font=\small\bfseries] at (\i,0) {\n}; \node[font=\small\bfseries,anchor=east] at (0.45,-\i) {\n};}
  \foreach \k in {0,...,#2} {\draw[gray] (\k+0.5,-0.5) -- (\k+0.5,-#2-0.5); \draw[gray] (0.5,-\k-0.5) -- (#2+0.5,-\k-0.5);}
  \ifloesung\ifstrempty{#1}{}{\foreach \z/\s/\t in {#1} {\node[text=red!75!black,font=\small] at (\s,-\z) {\t};}}\fi
\end{tikzpicture}}

% ---------- Java-Listings ----------
\lstset{language=Java,basicstyle=\ttfamily\small,keywordstyle=\bfseries,commentstyle=\color{gray!80!black},stringstyle=\color{black},showstringspaces=false,columns=flexible,keepspaces=true,tabsize=3,frame=single,framerule=0.4pt,rulecolor=\color{black},xleftmargin=2mm,framexleftmargin=1mm,aboveskip=2mm,belowskip=2mm,breaklines=true,escapechar=§,
 literate={ä}{{\"a}}1 {ö}{{\"o}}1 {ü}{{\"u}}1 {Ä}{{\"A}}1 {Ö}{{\"O}}1 {Ü}{{\"U}}1 {ß}{{\ss}}1 {–}{{--}}1 {„}{{\glqq}}1 {“}{{\grqq}}1}
\lstnewenvironment{java}[1][]{\lstset{#1}}{}

%  Übersetzen mit XeLaTeX, LuaLaTeX, pdfLaTeX oder Tectonic.
%
%  Blatt:   xelatex ab-3-1-grundbegriffe.tex
%  Lösung:  xelatex ab-3-1-grundbegriffe-lsg.tex
%           (Datei enthält nur: \def\mitloesung{}\documentclass[11pt,a4paper]{article}
\input{ab-vorlage}
% Dörfer-Graph (Merkkasten Begriffe): Koordinaten wie im Begriffe-Explorer der Website
\newcommand{\doerferknoten}{%
  \node[knoten,font=\tiny,minimum height=4mm,inner sep=1.5pt] (E) at (1.2,3.15) {Eulenhain};
  \node[knoten,font=\tiny,minimum height=4mm,inner sep=1.5pt] (S) at (3.3,3.25) {Sternwies};
  \node[knoten,font=\tiny,minimum height=4mm,inner sep=1.5pt] (K) at (5.2,2.65) {Kieselfurt};
  \node[knoten,font=\tiny,minimum height=4mm,inner sep=1.5pt] (F) at (2.5,1.85) {Funkenau};
  \node[knoten,font=\tiny,minimum height=4mm,inner sep=1.5pt] (H) at (4.8,1.25) {Haselbruck};
  \node[knoten,font=\tiny,minimum height=4mm,inner sep=1.5pt] (N) at (0.9,0.75) {Nebelried};
  \node[knoten,font=\tiny,minimum height=4mm,inner sep=1.5pt] (T) at (3.2,0.35) {Tannzapf};}
\begin{document}
\abtitel{3.1}{Graphen -- Grundbegriffe}

Hänsel und Gretel haben die Hexe überlistet. Doch die Brotkrumen, mit denen sie ihren Weg markiert hatten, haben die Vögel gefressen. Findest du trotzdem nach Hause?

\begin{aufgabe}[Hänsel und Gretel]{1}
Starte das Spiel auf der Website (Station 3.1).
\textbf{a)} Zeichne beim Spielen eine Karte des Waldes: jeder Ort ein Kästchen, jeder Weg eine Linie. Führt ein Weg nur in eine Richtung, zeichne einen Pfeil. Das Hexenhaus ist schon eingezeichnet.
\end{aufgabe}
\noindent\begin{tikzpicture}
\draw[step=5mm,gray!55,thin] (0,0) grid (17.4,7);
\node[knoten,fill=green!15] (hex) at (4,2.3) {Hexenhaus (Start)};
\ifloesung
  \node[lsgknoten] (warn) at (5.5,6.5) {Warnung};
  \node[lsgknoten] (wolf) at (10.5,6.5) {Werwolf};
  \node[lsgknoten] (stadt) at (1.6,5.2) {Stadt};
  \node[lsgknoten] (abz) at (5.5,5.2) {Abzweigung};
  \node[lsgknoten] (spinne) at (13,5.2) {Spinne};
  \node[lsgknoten] (baeck) at (1.6,3.6) {Bäckerei};
  \node[lsgknoten] (pilze) at (4,3.9) {Pilze};
  \node[lsgknoten] (w3) at (13,3.4) {Wald3};
  \node[lsgknoten] (w1) at (8,1.0) {Wald1};
  \node[lsgknoten] (stein) at (12.2,2.0) {Steinkreis};
  \node[lsgknoten] (w2) at (15,0.8) {Wald2};
  \node[lsgknoten] (heim) at (16.2,2.6) {Zuhause};
  \draw[lsgkante] (stadt) -- (baeck);
  \draw[lsgkante] (stadt) -- (pilze);
  \draw[lsgkante] (pilze) -- (abz);
  \draw[lsgkante] (pilze) -- (hex);
  \draw[lsggkante] (pilze) to[bend left=50] node[right,font=\scriptsize,text=red!75!black] {Pilze essen} (hex);
  \draw[lsgkante] (abz) -- (warn);
  \draw[lsggkante] (warn) -- (wolf);
  \draw[lsgkante] (abz) -- (spinne);
  \draw[lsgkante] (spinne) -- (w3);
  \draw[lsgkante] (hex) -- (w1);
  \draw[lsgkante] (w1) -- (w3);
  \draw[lsgkante] (w1) -- (w2);
  \draw[lsgkante] (w1) -- (stein);
  \draw[lsgkante] (w2) -- (w3);
  \draw[lsgkante] (w2) -- (stein);
  \draw[lsgkante] (w3) -- (stein);
  \draw[lsggkante] (w2) -- (heim);
  \node[text=red!75!black,font=\scriptsize,align=left,anchor=north west] at (0.2,1.3) {13 Knoten, 17 Kanten\\(14 ungerichtet, 3 gerichtet)};
\fi
\end{tikzpicture}

\textbf{b)} Am Ende zeigt dir das Spiel deinen Weg. Notiere ihn und die Anzahl der Schritte:
\lsglinien{1}{individuell; kürzester Weg: Hexenhaus -- Wald1 -- Wald2 -- Zuhause (3 Schritte)}

\begin{minipage}[t]{0.6\linewidth}
\begin{merke}[Graph und Kurzschreibweise $G = (V, E)$]\raggedright
Ein Graph besteht aus \lsg[2.3cm]{Knoten} (engl. \textit{vertices}) und \lsg[2.3cm]{Kanten} (engl. \textit{edges}). Eine Kante verbindet zwei Knoten. Im Spiel sind die Orte die \lsg[2cm]{Knoten}, die Wege die \lsg[2cm]{Kanten}.\\[1mm]
Kurzschreibweise $G = (V, E)$: $V$ ist die Menge der Knoten, $E$ die Menge der Kanten.
\begin{itemize}
\item $(A, B)$: \lsg[2.5cm]{gerichtete} Kante, nur von $A$ nach $B$
\item $\{A, B\}$: \lsg[2.5cm]{ungerichtete} Kante, steht für $(A, B)$ und \lsg[1.4cm]{$(B, A)$}
\end{itemize}
\end{merke}
\end{minipage}\hfill
\begin{minipage}[t]{0.37\linewidth}
\vspace{1mm}
\centering\small \textbf{Beispiel: Wer folgt wem?}\\[1mm]
\begin{tikzpicture}[x=0.95cm,y=0.9cm]
\node[knoten] (mira) at (2.4,2.5) {Mira};
\node[knoten] (ole) at (4.3,1.9) {Ole};
\node[knoten] (zeynep) at (3.4,0.3) {Zeynep};
\node[knoten] (finn) at (0.9,1.0) {Finn};
\node[knoten] (paula) at (0.5,2.5) {Paula};
\draw[dkante] (mira) -- (ole);
\draw[gkante] (ole) -- (zeynep);
\draw[gkante] (zeynep) -- (mira);
\draw[dkante] (mira) -- (finn);
\draw[gkante] (finn) -- (zeynep);
\end{tikzpicture}
\end{minipage}\\[1mm]
{\small $V = \{\text{Finn}, \text{Mira}, \text{Ole}, \text{Paula}, \text{Zeynep}\}$\\
$E = \{\{\text{Mira}, \text{Ole}\}, \{\text{Mira}, \text{Finn}\}, (\text{Ole}, \text{Zeynep}), (\text{Zeynep}, \text{Mira}), (\text{Finn}, \text{Zeynep})\}$}

\begin{aufgabe}[Hänsel und Gretel]{1}
\textbf{c)} Schreibe den Ausschnitt deiner Karte mit den Knoten Hexenhaus, Pilze, Stadt, Bäckerei und Wald1 in Kurzschreibweise auf (5 Kanten).\\[2mm]
$V = \{$\lsg[15.4cm]{Hexenhaus, Pilze, Stadt, Bäckerei, Wald1}$\}$\\[3mm]
$E = \{$\lsg[15.4cm]{\{Hexenhaus, Pilze\}, (Pilze, Hexenhaus), \{Pilze, Stadt\}, \{Stadt, Bäckerei\}, \{Hexenhaus, Wald1\}}$\}$
\end{aufgabe}

\newpage
\begin{minipage}[t]{0.71\linewidth}
\begin{merke}[Pfade, Zyklen, Richtung, Gewicht]\raggedright\small
\textbf{Pfad (Weg):} Folge von Knoten, in der aufeinanderfolgende Knoten durch eine \lsg[1.8cm]{Kante} verbunden sind; keine Kante wird doppelt benutzt.\\[0.6mm]
\textbf{Einfacher Pfad:} Pfad, in dem sich kein \lsg[1.8cm]{Knoten} wiederholt.\\
Beispiel in A: \lsg[5.6cm]{Nebelried -- Funkenau -- Kieselfurt -- Haselbruck}\\[0.6mm]
\textbf{Länge eines Pfades:} ungewichtet die \lsg[3.2cm]{Anzahl der Kanten}, gewichtet die \lsg[3.2cm]{Summe der Gewichte}.\\[0.6mm]
\textbf{Zyklus:} einfacher Pfad, dessen \lsg[5cm]{Start- und Endknoten gleich sind}.\\
Beispiel in A: \lsg[5.6cm]{Sternwies -- Funkenau -- Kieselfurt -- Sternwies}\\
Graph mit Zyklus: \emph{zyklisch}, sonst \lsg[3.4cm]{azyklisch (zyklenfrei)}.\\[0.6mm]
\textbf{Gerichtet:} Kanten nur \lsg[3cm]{in Pfeilrichtung} (Pfeil).\ \textbf{Ungerichtet:} \lsg[2.8cm]{in beide Richtungen} (Strich).\\[0.6mm]
\textbf{Stark zusammenhängend:} Von \lsg[2.4cm]{jedem Knoten} gibt es einen Pfad zu \lsg[3.2cm]{jedem anderen Knoten}.\\
\textbf{Schwach zusammenhängend:} \lsg[5.4cm]{nur ohne Beachtung der Richtungen}\\ \lsg[\linewidth]{hängt alles zusammen. B ist schwach: Von Tannzapf geht keine Kante weg.}\\[0.6mm]
\textbf{Gewichtet:} Jede Kante trägt eine \lsg[2.6cm]{Zahl (Gewicht)}, z.\,B. \lsg[3.6cm]{Entfernung, Fahrzeit}.
\end{merke}
\end{minipage}\hfill
\begin{minipage}[t]{0.27\linewidth}
\vspace{1mm}
\textbf{\small A} {\scriptsize ungerichtet}\\
\begin{tikzpicture}[x=0.8cm,y=0.85cm]
\doerferknoten
\draw[kante] (E)--(S); \draw[kante] (E)--(F); \draw[kante] (S)--(F); \draw[kante] (S)--(K);
\draw[kante] (F)--(K); \draw[kante] (F)--(N); \draw[kante] (K)--(H); \draw[kante] (H)--(F);
\draw[kante] (N)--(T); \draw[kante] (H)--(T);
\end{tikzpicture}\\[2mm]
\textbf{\small B} {\scriptsize gerichtet}\\
\begin{tikzpicture}[x=0.8cm,y=0.85cm]
\doerferknoten
\draw[gkante] (E)--(S); \draw[gkante] (E)--(F); \draw[gkante] (S)--(F); \draw[gkante] (S)--(K);
\draw[gkante] (F)--(K); \draw[gkante] (F)--(N); \draw[gkante] (K)--(H); \draw[gkante] (H)--(F);
\draw[gkante] (N)--(T); \draw[gkante] (H)--(T);
\end{tikzpicture}
\end{minipage}

\begin{aufgabe}{2}
{\raggedright Zeichne $G = (V, E)$ mit $V = \{A, B, C, D, E, F\}$ und\\ $E = \{(A, B), (A, C), (A, E), (B, C), (C, C), (C, D), (E, C), (F, D)\}$. Beantworte die Fragen.\par}\vspace{1mm}
\begin{minipage}[t]{0.42\linewidth}
\vspace{0pt}
\begin{tikzpicture}
\draw[step=5mm,gray!55,thin] (0,0) grid (7,3);
\ifloesung
  \foreach \n/\x/\y/\t in {a/0.6/1.5/A, b/2/2.55/B, e/2/0.45/E, c/3.4/1.5/C, d/4.9/1.5/D, f/6.4/1.5/F} {
    \node[kknoten,minimum size=6mm,draw=red!75!black,text=red!75!black] (\n) at (\x,\y) {\t};}
  \draw[lsggkante] (a)--(b); \draw[lsggkante] (a)--(c); \draw[lsggkante] (a)--(e);
  \draw[lsggkante] (b)--(c); \draw[lsggkante] (c)--(d); \draw[lsggkante] (e)--(c);
  \draw[lsggkante] (f)--(d);
  \draw[lsggkante] (c) to[out=60,in=120,looseness=6] (c);
\fi
\end{tikzpicture}
\end{minipage}\hfill
\begin{minipage}[t]{0.55\linewidth}
\vspace{0pt}\small\raggedright
Gerichtet? Woran erkennbar? \lsg[3.9cm]{ja, Paare ( ) statt \{ \}}\\[1.5mm]
Ein Zyklus: \lsg[5.7cm]{(C, C), eine Schlinge (Länge 1)}\\[1.5mm]
Stark/schwach zusammenhängend? \lsg[3.6cm]{schwach: von D}\\[1.5mm]
\lsg[0.98\linewidth]{kommt man nirgends hin.}\\[1.5mm]
Länge des längsten einfachen Pfades: \lsg[3.2cm]{3 (A--B--C--D)}
\end{minipage}
\end{aufgabe}

\begin{aufgabe}[Zurück in den Wald]{3}
Deine Karte aus Aufgabe 1: \textbf{a)} gerichtet?\ \textbf{b)} ein Zyklus?\ \textbf{c)} stark oder schwach zusammenhängend?\ \textbf{d)} „Der Weg zurück … scheint verschwunden“: Ist die Kante Hexenhaus -- Wald1 gerichtet? Prüfe im Spiel.
\lsglinien{2}{\small a) gemischt; gerichtet nur (Pilze, Hexenhaus), (Warnung, Werwolf), (Wald2, Zuhause)\quad b) Wald1 -- Wald2 -- Wald3 -- Wald1\\ c) schwach: Von Zuhause/Werwolf führt keine Kante weg.\quad d) nein: Von Wald2 aus führt ein Weg zurück.}
\end{aufgabe}

\begin{aufgabe}[Ordne die Graphen ein.]{4}
\vspace{-1mm}
\small\renewcommand{\arraystretch}{1.15}
\begin{tabularx}{\linewidth}{|>{\centering\arraybackslash}m{5.6cm}|>{\centering\arraybackslash}m{1.7cm}|>{\centering\arraybackslash}m{1.1cm}|>{\centering\arraybackslash}m{1.1cm}|>{\centering\arraybackslash}m{1.1cm}|>{\centering\arraybackslash}X|}\hline
\textbf{Graph} & \scriptsize\textbf{zusammen-\newline hängend?}\newline stark/schwach/nein & \scriptsize\textbf{ge-\newline richtet?} & \scriptsize\textbf{ge-\newline wichtet?} & \scriptsize\textbf{zyklisch?} & \scriptsize\textbf{Länge eines längsten einfachen Pfades}\\\hline
\rule{0pt}{8mm}\textbf{I}\ \begin{tikzpicture}[baseline=(k5.base),x=0.85cm,y=0.55cm]
\foreach \n/\x/\y in {1/0/1,2/1/1,3/2/1,4/3/1,5/1/0,6/3/0} {\node[kknoten,minimum size=5mm,inner sep=0pt,font=\tiny] (k\n) at (\x,\y) {K\n};}
\draw[kante] (k1)--(k2)--(k3)--(k4); \draw[kante] (k2)--(k5);
\end{tikzpicture}
& \lsgoder{\lsgfarbe nein}{} & \lsgoder{\lsgfarbe nein}{} & \lsgoder{\lsgfarbe nein}{} & \lsgoder{\lsgfarbe nein}{} & \lsgoder{\lsgfarbe 3: K1--K2--K3--K4}{}\\\hline
\rule{0pt}{5mm}\textbf{II}\ $V = \{A, B, C\}$\newline $E =\{(A, B), (B, A), (C, B), (C, C)\}$
& \lsgoder{\lsgfarbe schwach}{} & \lsgoder{\lsgfarbe ja}{} & \lsgoder{\lsgfarbe nein}{} & \lsgoder{\lsgfarbe ja}{} & \lsgoder{\lsgfarbe 2: C$\to$B$\to$A}{}\\\hline
\rule{0pt}{8mm}\textbf{III}\ \begin{tikzpicture}[baseline=(k3.base),x=0.9cm,y=0.6cm]
\foreach \n/\x/\y in {1/0/1,2/2/1,3/1/0,4/3/0} {\node[kknoten,minimum size=5mm,inner sep=0pt,font=\tiny] (k\n) at (\x,\y) {K\n};}
\draw[kante] (k1) -- node[gewicht] {4} (k2); \draw[kante] (k2) -- node[gewicht] {1} (k3);
\draw[kante] (k3) -- node[gewicht] {5} (k1); \draw[kante] (k3) -- node[gewicht] {2} (k4);
\end{tikzpicture}
& \lsgoder{\lsgfarbe stark}{} & \lsgoder{\lsgfarbe nein}{} & \lsgoder{\lsgfarbe ja}{} & \lsgoder{\lsgfarbe ja}{} & \lsgoder{\lsgfarbe 11: K2--K1--K3--K4}{}\\\hline
\end{tabularx}
\end{aufgabe}
\end{document}
)
%           (füllt alle Lücken rot aus; siehe \lsg, \lsglinien, \lsgoder)
% =====================================================================
\usepackage[a4paper,left=18mm,right=18mm,top=26mm,bottom=17mm,headheight=15mm,headsep=4mm,footskip=8mm]{geometry}
\usepackage{iftex}
\ifPDFTeX
  \usepackage[T1]{fontenc}
  \usepackage[utf8]{inputenc}
  \usepackage{helvet}
  \usepackage[scaled=0.86]{beramono}
  \renewcommand{\familydefault}{\sfdefault}
\else
  \usepackage{fontspec}
  \setmainfont{texgyreheros}[Extension=.otf,UprightFont=*-regular,BoldFont=*-bold,ItalicFont=*-italic,BoldItalicFont=*-bolditalic]
  \setmonofont{DejaVuSansMono}[Extension=.ttf,UprightFont=*,BoldFont=*-Bold,ItalicFont=*-Oblique,BoldItalicFont=*-BoldOblique,Scale=0.86]
\fi
\usepackage[ngerman,shorthands=off]{babel}
\usepackage{xcolor,graphicx,tabularx,array,enumitem,fancyhdr,tikz,listings,multicol,calc,etoolbox}
\usepackage[most]{tcolorbox}
\usetikzlibrary{arrows.meta,positioning,calc,decorations.pathreplacing,shapes.geometric}
\usepackage[hidelinks]{hyperref}
\setlength{\parindent}{0pt}
\setlength{\parskip}{3pt}
\setlist{nosep,leftmargin=*}

% ---------- Lösungsschalter ----------
\newif\ifloesung
\ifdefined\mitloesung\loesungtrue\fi
\newcommand{\lsgfarbe}{\color{red!75!black}}

% ---------- Kopf- und Fußzeile (einheitliche Form) ----------
\newcommand{\lehrkraft}{Hau}
\newcommand{\themenbereich}{3. Graphen}
\newcommand{\abnummer}{}
\pagestyle{fancy}\fancyhf{}
\renewcommand{\headrulewidth}{0.4pt}
\fancyhead[L]{\small Klasse 11\\Datum:}
\fancyhead[C]{\small Themenbereich:\\\themenbereich}
\fancyhead[R]{\small \lehrkraft\ifloesung\\{\lsgfarbe\bfseries Lösung}\fi}
\fancyfoot[C]{\scriptsize edu-mrh.de\,·\,Informatik 11\,·\,Blatt \abnummer\ifloesung\ (Lösung)\fi\,·\,Seite \thepage}
\newcommand{\abtitel}[2]{\renewcommand{\abnummer}{#1}\begin{center}\Large\bfseries #1\quad\underline{#2}\end{center}\vspace{-1mm}}

% ---------- Boxen ----------
\newenvironment{aufgabe}[2][]{\begin{tcolorbox}[enhanced,breakable,colback=white,colframe=black,boxrule=0.7pt,arc=3mm,left=2mm,right=2mm,top=1.5mm,bottom=1.5mm,before skip=2.5mm,after skip=2.5mm]\textbf{Aufgabe #2}\ifstrempty{#1}{}{ \textit{#1}}\par\vspace{0.6mm}}{\end{tcolorbox}}
\newtcolorbox{merke}[1][Merke]{enhanced,breakable,colback=green!5,colframe=green!40!black,boxrule=0.8pt,arc=2mm,left=2mm,right=2mm,top=1.5mm,bottom=1.5mm,before skip=2.5mm,after skip=2.5mm,title={#1},fonttitle=\bfseries,coltitle=white}
\newtcolorbox{hinweis}{enhanced,breakable,colback=yellow!12,colframe=orange!70!black,boxrule=0.6pt,arc=2mm,left=2mm,right=2mm,top=1mm,bottom=1mm,before skip=2mm,after skip=2mm}
\newtcolorbox{ausblick}[1][Ausblick]{enhanced,breakable,colback=teal!5,colframe=teal!60!black,boxrule=0.8pt,arc=2mm,left=2mm,right=2mm,top=1.5mm,bottom=1.5mm,before skip=2.5mm,after skip=2.5mm,title={#1},fonttitle=\bfseries,coltitle=white}
\newcommand{\web}[1]{{\small\color{gray}\textit{Website: #1}}}

% ---------- Lücken, Linien, Karos ----------
\newcommand{\luecke}[1][3.5cm]{\rule[-0.5ex]{#1}{0.4pt}}
\newcommand{\linien}[1]{\par\vspace{1.5mm}\foreach \i in {1,...,#1}{\noindent\rule[0pt]{\linewidth}{0.3pt}\par\vspace{5.5mm}}\vspace{-3mm}}
\newcommand{\kariert}[1]{\par\vspace{1.5mm}\noindent\begin{tikzpicture}\draw[step=5mm,gray!55,thin] (0,0) grid (\linewidth-1pt,#1*5mm);\end{tikzpicture}\par}
\newcommand{\karobox}[2]{\begin{tikzpicture}\draw[step=5mm,gray!55,thin] (0,0) grid (#1,#2);\end{tikzpicture}}
\newcommand{\code}[1]{\texttt{#1}}

% Lücke mit Lösung: \lsg[Mindestbreite]{Lösungstext}
\newlength{\lsgbreite}
\newcommand{\lsg}[2][3.5cm]{\ifloesung\settowidth{\lsgbreite}{\,#2\,}\ifdim\lsgbreite<#1\setlength{\lsgbreite}{#1}\fi\underline{\makebox[\lsgbreite][l]{\lsgfarbe\,#2}}\else\luecke[#1]\fi}
% Schreiblinien mit Lösung: \lsglinien{Anzahl}{Lösungstext} (gleiche Höhe wie \linien)
\newcommand{\lsglinien}[2]{\ifloesung\par\vspace{1.5mm}\noindent\begin{minipage}[t][#1\dimexpr 6.9mm\relax][t]{\linewidth}\lsgfarbe #2\end{minipage}\par\vspace{-1.5mm}\else\linien{#1}\fi}
% Beliebiger Inhalt: \lsgoder{nur in der Lösung}{nur im Blatt}
\newcommand{\lsgoder}[2]{\ifloesung\leavevmode#1\else#2\fi}

% ---------- Ergänzung (aus informatik/10/java/arbeitsblaetter/ab-vorlage.tex) ----------
% Bestehende Makros oben bleiben unverändert, damit ältere Blätter gleich aussehen.
%  \restzeile / \lsgrest{Text}   Lücke bis zum Zeilenende
%  \lsgtext{Text}                nur in der Lösung sichtbar (z. B. Tabellenzellen)
%  \lueckentext                  größerer Zeilenabstand, Flattersatz für Absätze mit Lücken
%  \schreibzeile                 Strut für Tabellenzellen zum Ausfüllen (8 mm hoch)
%  \alt{9}, \neu{7}              überschriebener alter Wert / neuer Wert (Dijkstra-Tabellen)
\usepackage{colortbl,xspace}
\usepackage[normalem]{ulem}
\newcommand{\lsgrest}[1]{\leavevmode\unskip\hspace{1.5mm}\ifloesung\rlap{\raisebox{0.8pt}{\lsgfarbe\,#1}}\fi\leaders\hbox{\rule[-0.6ex]{0.5pt}{0.4pt}}\hfill\null\par}
\newcommand{\restzeile}{\lsgrest{}}
\newcommand{\lsgtext}[1]{\ifloesung\leavevmode{\lsgfarbe#1}\fi}
\newcommand{\lueckentext}{\linespread{1.5}\selectfont\raggedright}
\newcommand{\schreibzeile}{\rule[-2.8mm]{0pt}{8mm}}
\newcommand{\alt}[1]{\leavevmode{\color{blue!60!black}\sout{#1}}}
\newcommand{\neu}[1]{\leavevmode{\color{green!50!black}\bfseries #1}}
\colorlet{djaktuell}{orange!40}
\colorlet{djfertig}{blue!22}

% ---------- Graphen zeichnen ----------
\tikzset{
  knoten/.style={draw,rounded corners=3pt,fill=orange!12,minimum height=6mm,inner sep=3pt,font=\small},
  kknoten/.style={draw,circle,fill=orange!12,minimum size=7mm,inner sep=1pt,font=\small},
  kante/.style={thick},
  gkante/.style={thick,-{Stealth[length=2.6mm]}},
  dkante/.style={thick,{Stealth[length=2.6mm]}-{Stealth[length=2.6mm]}},
  gewicht/.style={font=\scriptsize,fill=white,inner sep=1pt,text=teal!60!black},
  lsgkante/.style={thick,red!75!black},
  lsggkante/.style={thick,red!75!black,-{Stealth[length=2.6mm]}},
  lsgknoten/.style={knoten,draw=red!75!black,text=red!75!black,fill=red!4},
}
% Adjazenzmatrix zum Ausfüllen: \adjmx[Lösungseinträge]{Anzahl}{Knotenliste}{Zellbreite}
% Lösungseinträge als Zeile/Spalte/Text, z. B. \adjmx[1/2/X,2/1/X]{3}{A,B,C}{7mm}
\newcommand{\adjmx}[4][]{\begin{tikzpicture}[x=#4,y=#4]
  \foreach \n [count=\i] in {#3} {\node[font=\small\bfseries] at (\i,0) {\n}; \node[font=\small\bfseries,anchor=east] at (0.45,-\i) {\n};}
  \foreach \k in {0,...,#2} {\draw[gray] (\k+0.5,-0.5) -- (\k+0.5,-#2-0.5); \draw[gray] (0.5,-\k-0.5) -- (#2+0.5,-\k-0.5);}
  \ifloesung\ifstrempty{#1}{}{\foreach \z/\s/\t in {#1} {\node[text=red!75!black,font=\small] at (\s,-\z) {\t};}}\fi
\end{tikzpicture}}

% ---------- Java-Listings ----------
\lstset{language=Java,basicstyle=\ttfamily\small,keywordstyle=\bfseries,commentstyle=\color{gray!80!black},stringstyle=\color{black},showstringspaces=false,columns=flexible,keepspaces=true,tabsize=3,frame=single,framerule=0.4pt,rulecolor=\color{black},xleftmargin=2mm,framexleftmargin=1mm,aboveskip=2mm,belowskip=2mm,breaklines=true,escapechar=§,
 literate={ä}{{\"a}}1 {ö}{{\"o}}1 {ü}{{\"u}}1 {Ä}{{\"A}}1 {Ö}{{\"O}}1 {Ü}{{\"U}}1 {ß}{{\ss}}1 {–}{{--}}1 {„}{{\glqq}}1 {“}{{\grqq}}1}
\lstnewenvironment{java}[1][]{\lstset{#1}}{}

%  Übersetzen mit XeLaTeX, LuaLaTeX, pdfLaTeX oder Tectonic.
%
%  Blatt:   xelatex ab-3-1-grundbegriffe.tex
%  Lösung:  xelatex ab-3-1-grundbegriffe-lsg.tex
%           (Datei enthält nur: \def\mitloesung{}\documentclass[11pt,a4paper]{article}
% =====================================================================
%  Gemeinsame Vorlage der Arbeitsblätter "Graphen" – Informatik 11 (NTG)
%  (gleiches Layout wie informatik/10/java/arbeitsblaetter/ab-vorlage.tex)
%  Einbinden in jedem Blatt direkt nach \documentclass: \input{ab-vorlage}
%  Übersetzen mit XeLaTeX, LuaLaTeX, pdfLaTeX oder Tectonic.
%
%  Blatt:   xelatex ab-3-1-grundbegriffe.tex
%  Lösung:  xelatex ab-3-1-grundbegriffe-lsg.tex
%           (Datei enthält nur: \def\mitloesung{}\input{ab-3-1-grundbegriffe})
%           (füllt alle Lücken rot aus; siehe \lsg, \lsglinien, \lsgoder)
% =====================================================================
\usepackage[a4paper,left=18mm,right=18mm,top=26mm,bottom=17mm,headheight=15mm,headsep=4mm,footskip=8mm]{geometry}
\usepackage{iftex}
\ifPDFTeX
  \usepackage[T1]{fontenc}
  \usepackage[utf8]{inputenc}
  \usepackage{helvet}
  \usepackage[scaled=0.86]{beramono}
  \renewcommand{\familydefault}{\sfdefault}
\else
  \usepackage{fontspec}
  \setmainfont{texgyreheros}[Extension=.otf,UprightFont=*-regular,BoldFont=*-bold,ItalicFont=*-italic,BoldItalicFont=*-bolditalic]
  \setmonofont{DejaVuSansMono}[Extension=.ttf,UprightFont=*,BoldFont=*-Bold,ItalicFont=*-Oblique,BoldItalicFont=*-BoldOblique,Scale=0.86]
\fi
\usepackage[ngerman,shorthands=off]{babel}
\usepackage{xcolor,graphicx,tabularx,array,enumitem,fancyhdr,tikz,listings,multicol,calc,etoolbox}
\usepackage[most]{tcolorbox}
\usetikzlibrary{arrows.meta,positioning,calc,decorations.pathreplacing,shapes.geometric}
\usepackage[hidelinks]{hyperref}
\setlength{\parindent}{0pt}
\setlength{\parskip}{3pt}
\setlist{nosep,leftmargin=*}

% ---------- Lösungsschalter ----------
\newif\ifloesung
\ifdefined\mitloesung\loesungtrue\fi
\newcommand{\lsgfarbe}{\color{red!75!black}}

% ---------- Kopf- und Fußzeile (einheitliche Form) ----------
\newcommand{\lehrkraft}{Hau}
\newcommand{\themenbereich}{3. Graphen}
\newcommand{\abnummer}{}
\pagestyle{fancy}\fancyhf{}
\renewcommand{\headrulewidth}{0.4pt}
\fancyhead[L]{\small Klasse 11\\Datum:}
\fancyhead[C]{\small Themenbereich:\\\themenbereich}
\fancyhead[R]{\small \lehrkraft\ifloesung\\{\lsgfarbe\bfseries Lösung}\fi}
\fancyfoot[C]{\scriptsize edu-mrh.de\,·\,Informatik 11\,·\,Blatt \abnummer\ifloesung\ (Lösung)\fi\,·\,Seite \thepage}
\newcommand{\abtitel}[2]{\renewcommand{\abnummer}{#1}\begin{center}\Large\bfseries #1\quad\underline{#2}\end{center}\vspace{-1mm}}

% ---------- Boxen ----------
\newenvironment{aufgabe}[2][]{\begin{tcolorbox}[enhanced,breakable,colback=white,colframe=black,boxrule=0.7pt,arc=3mm,left=2mm,right=2mm,top=1.5mm,bottom=1.5mm,before skip=2.5mm,after skip=2.5mm]\textbf{Aufgabe #2}\ifstrempty{#1}{}{ \textit{#1}}\par\vspace{0.6mm}}{\end{tcolorbox}}
\newtcolorbox{merke}[1][Merke]{enhanced,breakable,colback=green!5,colframe=green!40!black,boxrule=0.8pt,arc=2mm,left=2mm,right=2mm,top=1.5mm,bottom=1.5mm,before skip=2.5mm,after skip=2.5mm,title={#1},fonttitle=\bfseries,coltitle=white}
\newtcolorbox{hinweis}{enhanced,breakable,colback=yellow!12,colframe=orange!70!black,boxrule=0.6pt,arc=2mm,left=2mm,right=2mm,top=1mm,bottom=1mm,before skip=2mm,after skip=2mm}
\newtcolorbox{ausblick}[1][Ausblick]{enhanced,breakable,colback=teal!5,colframe=teal!60!black,boxrule=0.8pt,arc=2mm,left=2mm,right=2mm,top=1.5mm,bottom=1.5mm,before skip=2.5mm,after skip=2.5mm,title={#1},fonttitle=\bfseries,coltitle=white}
\newcommand{\web}[1]{{\small\color{gray}\textit{Website: #1}}}

% ---------- Lücken, Linien, Karos ----------
\newcommand{\luecke}[1][3.5cm]{\rule[-0.5ex]{#1}{0.4pt}}
\newcommand{\linien}[1]{\par\vspace{1.5mm}\foreach \i in {1,...,#1}{\noindent\rule[0pt]{\linewidth}{0.3pt}\par\vspace{5.5mm}}\vspace{-3mm}}
\newcommand{\kariert}[1]{\par\vspace{1.5mm}\noindent\begin{tikzpicture}\draw[step=5mm,gray!55,thin] (0,0) grid (\linewidth-1pt,#1*5mm);\end{tikzpicture}\par}
\newcommand{\karobox}[2]{\begin{tikzpicture}\draw[step=5mm,gray!55,thin] (0,0) grid (#1,#2);\end{tikzpicture}}
\newcommand{\code}[1]{\texttt{#1}}

% Lücke mit Lösung: \lsg[Mindestbreite]{Lösungstext}
\newlength{\lsgbreite}
\newcommand{\lsg}[2][3.5cm]{\ifloesung\settowidth{\lsgbreite}{\,#2\,}\ifdim\lsgbreite<#1\setlength{\lsgbreite}{#1}\fi\underline{\makebox[\lsgbreite][l]{\lsgfarbe\,#2}}\else\luecke[#1]\fi}
% Schreiblinien mit Lösung: \lsglinien{Anzahl}{Lösungstext} (gleiche Höhe wie \linien)
\newcommand{\lsglinien}[2]{\ifloesung\par\vspace{1.5mm}\noindent\begin{minipage}[t][#1\dimexpr 6.9mm\relax][t]{\linewidth}\lsgfarbe #2\end{minipage}\par\vspace{-1.5mm}\else\linien{#1}\fi}
% Beliebiger Inhalt: \lsgoder{nur in der Lösung}{nur im Blatt}
\newcommand{\lsgoder}[2]{\ifloesung\leavevmode#1\else#2\fi}

% ---------- Ergänzung (aus informatik/10/java/arbeitsblaetter/ab-vorlage.tex) ----------
% Bestehende Makros oben bleiben unverändert, damit ältere Blätter gleich aussehen.
%  \restzeile / \lsgrest{Text}   Lücke bis zum Zeilenende
%  \lsgtext{Text}                nur in der Lösung sichtbar (z. B. Tabellenzellen)
%  \lueckentext                  größerer Zeilenabstand, Flattersatz für Absätze mit Lücken
%  \schreibzeile                 Strut für Tabellenzellen zum Ausfüllen (8 mm hoch)
%  \alt{9}, \neu{7}              überschriebener alter Wert / neuer Wert (Dijkstra-Tabellen)
\usepackage{colortbl,xspace}
\usepackage[normalem]{ulem}
\newcommand{\lsgrest}[1]{\leavevmode\unskip\hspace{1.5mm}\ifloesung\rlap{\raisebox{0.8pt}{\lsgfarbe\,#1}}\fi\leaders\hbox{\rule[-0.6ex]{0.5pt}{0.4pt}}\hfill\null\par}
\newcommand{\restzeile}{\lsgrest{}}
\newcommand{\lsgtext}[1]{\ifloesung\leavevmode{\lsgfarbe#1}\fi}
\newcommand{\lueckentext}{\linespread{1.5}\selectfont\raggedright}
\newcommand{\schreibzeile}{\rule[-2.8mm]{0pt}{8mm}}
\newcommand{\alt}[1]{\leavevmode{\color{blue!60!black}\sout{#1}}}
\newcommand{\neu}[1]{\leavevmode{\color{green!50!black}\bfseries #1}}
\colorlet{djaktuell}{orange!40}
\colorlet{djfertig}{blue!22}

% ---------- Graphen zeichnen ----------
\tikzset{
  knoten/.style={draw,rounded corners=3pt,fill=orange!12,minimum height=6mm,inner sep=3pt,font=\small},
  kknoten/.style={draw,circle,fill=orange!12,minimum size=7mm,inner sep=1pt,font=\small},
  kante/.style={thick},
  gkante/.style={thick,-{Stealth[length=2.6mm]}},
  dkante/.style={thick,{Stealth[length=2.6mm]}-{Stealth[length=2.6mm]}},
  gewicht/.style={font=\scriptsize,fill=white,inner sep=1pt,text=teal!60!black},
  lsgkante/.style={thick,red!75!black},
  lsggkante/.style={thick,red!75!black,-{Stealth[length=2.6mm]}},
  lsgknoten/.style={knoten,draw=red!75!black,text=red!75!black,fill=red!4},
}
% Adjazenzmatrix zum Ausfüllen: \adjmx[Lösungseinträge]{Anzahl}{Knotenliste}{Zellbreite}
% Lösungseinträge als Zeile/Spalte/Text, z. B. \adjmx[1/2/X,2/1/X]{3}{A,B,C}{7mm}
\newcommand{\adjmx}[4][]{\begin{tikzpicture}[x=#4,y=#4]
  \foreach \n [count=\i] in {#3} {\node[font=\small\bfseries] at (\i,0) {\n}; \node[font=\small\bfseries,anchor=east] at (0.45,-\i) {\n};}
  \foreach \k in {0,...,#2} {\draw[gray] (\k+0.5,-0.5) -- (\k+0.5,-#2-0.5); \draw[gray] (0.5,-\k-0.5) -- (#2+0.5,-\k-0.5);}
  \ifloesung\ifstrempty{#1}{}{\foreach \z/\s/\t in {#1} {\node[text=red!75!black,font=\small] at (\s,-\z) {\t};}}\fi
\end{tikzpicture}}

% ---------- Java-Listings ----------
\lstset{language=Java,basicstyle=\ttfamily\small,keywordstyle=\bfseries,commentstyle=\color{gray!80!black},stringstyle=\color{black},showstringspaces=false,columns=flexible,keepspaces=true,tabsize=3,frame=single,framerule=0.4pt,rulecolor=\color{black},xleftmargin=2mm,framexleftmargin=1mm,aboveskip=2mm,belowskip=2mm,breaklines=true,escapechar=§,
 literate={ä}{{\"a}}1 {ö}{{\"o}}1 {ü}{{\"u}}1 {Ä}{{\"A}}1 {Ö}{{\"O}}1 {Ü}{{\"U}}1 {ß}{{\ss}}1 {–}{{--}}1 {„}{{\glqq}}1 {“}{{\grqq}}1}
\lstnewenvironment{java}[1][]{\lstset{#1}}{}

% Dörfer-Graph (Merkkasten Begriffe): Koordinaten wie im Begriffe-Explorer der Website
\newcommand{\doerferknoten}{%
  \node[knoten,font=\tiny,minimum height=4mm,inner sep=1.5pt] (E) at (1.2,3.15) {Eulenhain};
  \node[knoten,font=\tiny,minimum height=4mm,inner sep=1.5pt] (S) at (3.3,3.25) {Sternwies};
  \node[knoten,font=\tiny,minimum height=4mm,inner sep=1.5pt] (K) at (5.2,2.65) {Kieselfurt};
  \node[knoten,font=\tiny,minimum height=4mm,inner sep=1.5pt] (F) at (2.5,1.85) {Funkenau};
  \node[knoten,font=\tiny,minimum height=4mm,inner sep=1.5pt] (H) at (4.8,1.25) {Haselbruck};
  \node[knoten,font=\tiny,minimum height=4mm,inner sep=1.5pt] (N) at (0.9,0.75) {Nebelried};
  \node[knoten,font=\tiny,minimum height=4mm,inner sep=1.5pt] (T) at (3.2,0.35) {Tannzapf};}
\begin{document}
\abtitel{3.1}{Graphen -- Grundbegriffe}

Hänsel und Gretel haben die Hexe überlistet. Doch die Brotkrumen, mit denen sie ihren Weg markiert hatten, haben die Vögel gefressen. Findest du trotzdem nach Hause?

\begin{aufgabe}[Hänsel und Gretel]{1}
Starte das Spiel auf der Website (Station 3.1).
\textbf{a)} Zeichne beim Spielen eine Karte des Waldes: jeder Ort ein Kästchen, jeder Weg eine Linie. Führt ein Weg nur in eine Richtung, zeichne einen Pfeil. Das Hexenhaus ist schon eingezeichnet.
\end{aufgabe}
\noindent\begin{tikzpicture}
\draw[step=5mm,gray!55,thin] (0,0) grid (17.4,7);
\node[knoten,fill=green!15] (hex) at (4,2.3) {Hexenhaus (Start)};
\ifloesung
  \node[lsgknoten] (warn) at (5.5,6.5) {Warnung};
  \node[lsgknoten] (wolf) at (10.5,6.5) {Werwolf};
  \node[lsgknoten] (stadt) at (1.6,5.2) {Stadt};
  \node[lsgknoten] (abz) at (5.5,5.2) {Abzweigung};
  \node[lsgknoten] (spinne) at (13,5.2) {Spinne};
  \node[lsgknoten] (baeck) at (1.6,3.6) {Bäckerei};
  \node[lsgknoten] (pilze) at (4,3.9) {Pilze};
  \node[lsgknoten] (w3) at (13,3.4) {Wald3};
  \node[lsgknoten] (w1) at (8,1.0) {Wald1};
  \node[lsgknoten] (stein) at (12.2,2.0) {Steinkreis};
  \node[lsgknoten] (w2) at (15,0.8) {Wald2};
  \node[lsgknoten] (heim) at (16.2,2.6) {Zuhause};
  \draw[lsgkante] (stadt) -- (baeck);
  \draw[lsgkante] (stadt) -- (pilze);
  \draw[lsgkante] (pilze) -- (abz);
  \draw[lsgkante] (pilze) -- (hex);
  \draw[lsggkante] (pilze) to[bend left=50] node[right,font=\scriptsize,text=red!75!black] {Pilze essen} (hex);
  \draw[lsgkante] (abz) -- (warn);
  \draw[lsggkante] (warn) -- (wolf);
  \draw[lsgkante] (abz) -- (spinne);
  \draw[lsgkante] (spinne) -- (w3);
  \draw[lsgkante] (hex) -- (w1);
  \draw[lsgkante] (w1) -- (w3);
  \draw[lsgkante] (w1) -- (w2);
  \draw[lsgkante] (w1) -- (stein);
  \draw[lsgkante] (w2) -- (w3);
  \draw[lsgkante] (w2) -- (stein);
  \draw[lsgkante] (w3) -- (stein);
  \draw[lsggkante] (w2) -- (heim);
  \node[text=red!75!black,font=\scriptsize,align=left,anchor=north west] at (0.2,1.3) {13 Knoten, 17 Kanten\\(14 ungerichtet, 3 gerichtet)};
\fi
\end{tikzpicture}

\textbf{b)} Am Ende zeigt dir das Spiel deinen Weg. Notiere ihn und die Anzahl der Schritte:
\lsglinien{1}{individuell; kürzester Weg: Hexenhaus -- Wald1 -- Wald2 -- Zuhause (3 Schritte)}

\begin{minipage}[t]{0.6\linewidth}
\begin{merke}[Graph und Kurzschreibweise $G = (V, E)$]\raggedright
Ein Graph besteht aus \lsg[2.3cm]{Knoten} (engl. \textit{vertices}) und \lsg[2.3cm]{Kanten} (engl. \textit{edges}). Eine Kante verbindet zwei Knoten. Im Spiel sind die Orte die \lsg[2cm]{Knoten}, die Wege die \lsg[2cm]{Kanten}.\\[1mm]
Kurzschreibweise $G = (V, E)$: $V$ ist die Menge der Knoten, $E$ die Menge der Kanten.
\begin{itemize}
\item $(A, B)$: \lsg[2.5cm]{gerichtete} Kante, nur von $A$ nach $B$
\item $\{A, B\}$: \lsg[2.5cm]{ungerichtete} Kante, steht für $(A, B)$ und \lsg[1.4cm]{$(B, A)$}
\end{itemize}
\end{merke}
\end{minipage}\hfill
\begin{minipage}[t]{0.37\linewidth}
\vspace{1mm}
\centering\small \textbf{Beispiel: Wer folgt wem?}\\[1mm]
\begin{tikzpicture}[x=0.95cm,y=0.9cm]
\node[knoten] (mira) at (2.4,2.5) {Mira};
\node[knoten] (ole) at (4.3,1.9) {Ole};
\node[knoten] (zeynep) at (3.4,0.3) {Zeynep};
\node[knoten] (finn) at (0.9,1.0) {Finn};
\node[knoten] (paula) at (0.5,2.5) {Paula};
\draw[dkante] (mira) -- (ole);
\draw[gkante] (ole) -- (zeynep);
\draw[gkante] (zeynep) -- (mira);
\draw[dkante] (mira) -- (finn);
\draw[gkante] (finn) -- (zeynep);
\end{tikzpicture}
\end{minipage}\\[1mm]
{\small $V = \{\text{Finn}, \text{Mira}, \text{Ole}, \text{Paula}, \text{Zeynep}\}$\\
$E = \{\{\text{Mira}, \text{Ole}\}, \{\text{Mira}, \text{Finn}\}, (\text{Ole}, \text{Zeynep}), (\text{Zeynep}, \text{Mira}), (\text{Finn}, \text{Zeynep})\}$}

\begin{aufgabe}[Hänsel und Gretel]{1}
\textbf{c)} Schreibe den Ausschnitt deiner Karte mit den Knoten Hexenhaus, Pilze, Stadt, Bäckerei und Wald1 in Kurzschreibweise auf (5 Kanten).\\[2mm]
$V = \{$\lsg[15.4cm]{Hexenhaus, Pilze, Stadt, Bäckerei, Wald1}$\}$\\[3mm]
$E = \{$\lsg[15.4cm]{\{Hexenhaus, Pilze\}, (Pilze, Hexenhaus), \{Pilze, Stadt\}, \{Stadt, Bäckerei\}, \{Hexenhaus, Wald1\}}$\}$
\end{aufgabe}

\newpage
\begin{minipage}[t]{0.71\linewidth}
\begin{merke}[Pfade, Zyklen, Richtung, Gewicht]\raggedright\small
\textbf{Pfad (Weg):} Folge von Knoten, in der aufeinanderfolgende Knoten durch eine \lsg[1.8cm]{Kante} verbunden sind; keine Kante wird doppelt benutzt.\\[0.6mm]
\textbf{Einfacher Pfad:} Pfad, in dem sich kein \lsg[1.8cm]{Knoten} wiederholt.\\
Beispiel in A: \lsg[5.6cm]{Nebelried -- Funkenau -- Kieselfurt -- Haselbruck}\\[0.6mm]
\textbf{Länge eines Pfades:} ungewichtet die \lsg[3.2cm]{Anzahl der Kanten}, gewichtet die \lsg[3.2cm]{Summe der Gewichte}.\\[0.6mm]
\textbf{Zyklus:} einfacher Pfad, dessen \lsg[5cm]{Start- und Endknoten gleich sind}.\\
Beispiel in A: \lsg[5.6cm]{Sternwies -- Funkenau -- Kieselfurt -- Sternwies}\\
Graph mit Zyklus: \emph{zyklisch}, sonst \lsg[3.4cm]{azyklisch (zyklenfrei)}.\\[0.6mm]
\textbf{Gerichtet:} Kanten nur \lsg[3cm]{in Pfeilrichtung} (Pfeil).\ \textbf{Ungerichtet:} \lsg[2.8cm]{in beide Richtungen} (Strich).\\[0.6mm]
\textbf{Stark zusammenhängend:} Von \lsg[2.4cm]{jedem Knoten} gibt es einen Pfad zu \lsg[3.2cm]{jedem anderen Knoten}.\\
\textbf{Schwach zusammenhängend:} \lsg[5.4cm]{nur ohne Beachtung der Richtungen}\\ \lsg[\linewidth]{hängt alles zusammen. B ist schwach: Von Tannzapf geht keine Kante weg.}\\[0.6mm]
\textbf{Gewichtet:} Jede Kante trägt eine \lsg[2.6cm]{Zahl (Gewicht)}, z.\,B. \lsg[3.6cm]{Entfernung, Fahrzeit}.
\end{merke}
\end{minipage}\hfill
\begin{minipage}[t]{0.27\linewidth}
\vspace{1mm}
\textbf{\small A} {\scriptsize ungerichtet}\\
\begin{tikzpicture}[x=0.8cm,y=0.85cm]
\doerferknoten
\draw[kante] (E)--(S); \draw[kante] (E)--(F); \draw[kante] (S)--(F); \draw[kante] (S)--(K);
\draw[kante] (F)--(K); \draw[kante] (F)--(N); \draw[kante] (K)--(H); \draw[kante] (H)--(F);
\draw[kante] (N)--(T); \draw[kante] (H)--(T);
\end{tikzpicture}\\[2mm]
\textbf{\small B} {\scriptsize gerichtet}\\
\begin{tikzpicture}[x=0.8cm,y=0.85cm]
\doerferknoten
\draw[gkante] (E)--(S); \draw[gkante] (E)--(F); \draw[gkante] (S)--(F); \draw[gkante] (S)--(K);
\draw[gkante] (F)--(K); \draw[gkante] (F)--(N); \draw[gkante] (K)--(H); \draw[gkante] (H)--(F);
\draw[gkante] (N)--(T); \draw[gkante] (H)--(T);
\end{tikzpicture}
\end{minipage}

\begin{aufgabe}{2}
{\raggedright Zeichne $G = (V, E)$ mit $V = \{A, B, C, D, E, F\}$ und\\ $E = \{(A, B), (A, C), (A, E), (B, C), (C, C), (C, D), (E, C), (F, D)\}$. Beantworte die Fragen.\par}\vspace{1mm}
\begin{minipage}[t]{0.42\linewidth}
\vspace{0pt}
\begin{tikzpicture}
\draw[step=5mm,gray!55,thin] (0,0) grid (7,3);
\ifloesung
  \foreach \n/\x/\y/\t in {a/0.6/1.5/A, b/2/2.55/B, e/2/0.45/E, c/3.4/1.5/C, d/4.9/1.5/D, f/6.4/1.5/F} {
    \node[kknoten,minimum size=6mm,draw=red!75!black,text=red!75!black] (\n) at (\x,\y) {\t};}
  \draw[lsggkante] (a)--(b); \draw[lsggkante] (a)--(c); \draw[lsggkante] (a)--(e);
  \draw[lsggkante] (b)--(c); \draw[lsggkante] (c)--(d); \draw[lsggkante] (e)--(c);
  \draw[lsggkante] (f)--(d);
  \draw[lsggkante] (c) to[out=60,in=120,looseness=6] (c);
\fi
\end{tikzpicture}
\end{minipage}\hfill
\begin{minipage}[t]{0.55\linewidth}
\vspace{0pt}\small\raggedright
Gerichtet? Woran erkennbar? \lsg[3.9cm]{ja, Paare ( ) statt \{ \}}\\[1.5mm]
Ein Zyklus: \lsg[5.7cm]{(C, C), eine Schlinge (Länge 1)}\\[1.5mm]
Stark/schwach zusammenhängend? \lsg[3.6cm]{schwach: von D}\\[1.5mm]
\lsg[0.98\linewidth]{kommt man nirgends hin.}\\[1.5mm]
Länge des längsten einfachen Pfades: \lsg[3.2cm]{3 (A--B--C--D)}
\end{minipage}
\end{aufgabe}

\begin{aufgabe}[Zurück in den Wald]{3}
Deine Karte aus Aufgabe 1: \textbf{a)} gerichtet?\ \textbf{b)} ein Zyklus?\ \textbf{c)} stark oder schwach zusammenhängend?\ \textbf{d)} „Der Weg zurück … scheint verschwunden“: Ist die Kante Hexenhaus -- Wald1 gerichtet? Prüfe im Spiel.
\lsglinien{2}{\small a) gemischt; gerichtet nur (Pilze, Hexenhaus), (Warnung, Werwolf), (Wald2, Zuhause)\quad b) Wald1 -- Wald2 -- Wald3 -- Wald1\\ c) schwach: Von Zuhause/Werwolf führt keine Kante weg.\quad d) nein: Von Wald2 aus führt ein Weg zurück.}
\end{aufgabe}

\begin{aufgabe}[Ordne die Graphen ein.]{4}
\vspace{-1mm}
\small\renewcommand{\arraystretch}{1.15}
\begin{tabularx}{\linewidth}{|>{\centering\arraybackslash}m{5.6cm}|>{\centering\arraybackslash}m{1.7cm}|>{\centering\arraybackslash}m{1.1cm}|>{\centering\arraybackslash}m{1.1cm}|>{\centering\arraybackslash}m{1.1cm}|>{\centering\arraybackslash}X|}\hline
\textbf{Graph} & \scriptsize\textbf{zusammen-\newline hängend?}\newline stark/schwach/nein & \scriptsize\textbf{ge-\newline richtet?} & \scriptsize\textbf{ge-\newline wichtet?} & \scriptsize\textbf{zyklisch?} & \scriptsize\textbf{Länge eines längsten einfachen Pfades}\\\hline
\rule{0pt}{8mm}\textbf{I}\ \begin{tikzpicture}[baseline=(k5.base),x=0.85cm,y=0.55cm]
\foreach \n/\x/\y in {1/0/1,2/1/1,3/2/1,4/3/1,5/1/0,6/3/0} {\node[kknoten,minimum size=5mm,inner sep=0pt,font=\tiny] (k\n) at (\x,\y) {K\n};}
\draw[kante] (k1)--(k2)--(k3)--(k4); \draw[kante] (k2)--(k5);
\end{tikzpicture}
& \lsgoder{\lsgfarbe nein}{} & \lsgoder{\lsgfarbe nein}{} & \lsgoder{\lsgfarbe nein}{} & \lsgoder{\lsgfarbe nein}{} & \lsgoder{\lsgfarbe 3: K1--K2--K3--K4}{}\\\hline
\rule{0pt}{5mm}\textbf{II}\ $V = \{A, B, C\}$\newline $E =\{(A, B), (B, A), (C, B), (C, C)\}$
& \lsgoder{\lsgfarbe schwach}{} & \lsgoder{\lsgfarbe ja}{} & \lsgoder{\lsgfarbe nein}{} & \lsgoder{\lsgfarbe ja}{} & \lsgoder{\lsgfarbe 2: C$\to$B$\to$A}{}\\\hline
\rule{0pt}{8mm}\textbf{III}\ \begin{tikzpicture}[baseline=(k3.base),x=0.9cm,y=0.6cm]
\foreach \n/\x/\y in {1/0/1,2/2/1,3/1/0,4/3/0} {\node[kknoten,minimum size=5mm,inner sep=0pt,font=\tiny] (k\n) at (\x,\y) {K\n};}
\draw[kante] (k1) -- node[gewicht] {4} (k2); \draw[kante] (k2) -- node[gewicht] {1} (k3);
\draw[kante] (k3) -- node[gewicht] {5} (k1); \draw[kante] (k3) -- node[gewicht] {2} (k4);
\end{tikzpicture}
& \lsgoder{\lsgfarbe stark}{} & \lsgoder{\lsgfarbe nein}{} & \lsgoder{\lsgfarbe ja}{} & \lsgoder{\lsgfarbe ja}{} & \lsgoder{\lsgfarbe 11: K2--K1--K3--K4}{}\\\hline
\end{tabularx}
\end{aufgabe}
\end{document}
)
%           (füllt alle Lücken rot aus; siehe \lsg, \lsglinien, \lsgoder)
% =====================================================================
\usepackage[a4paper,left=18mm,right=18mm,top=26mm,bottom=17mm,headheight=15mm,headsep=4mm,footskip=8mm]{geometry}
\usepackage{iftex}
\ifPDFTeX
  \usepackage[T1]{fontenc}
  \usepackage[utf8]{inputenc}
  \usepackage{helvet}
  \usepackage[scaled=0.86]{beramono}
  \renewcommand{\familydefault}{\sfdefault}
\else
  \usepackage{fontspec}
  \setmainfont{texgyreheros}[Extension=.otf,UprightFont=*-regular,BoldFont=*-bold,ItalicFont=*-italic,BoldItalicFont=*-bolditalic]
  \setmonofont{DejaVuSansMono}[Extension=.ttf,UprightFont=*,BoldFont=*-Bold,ItalicFont=*-Oblique,BoldItalicFont=*-BoldOblique,Scale=0.86]
\fi
\usepackage[ngerman,shorthands=off]{babel}
\usepackage{xcolor,graphicx,tabularx,array,enumitem,fancyhdr,tikz,listings,multicol,calc,etoolbox}
\usepackage[most]{tcolorbox}
\usetikzlibrary{arrows.meta,positioning,calc,decorations.pathreplacing,shapes.geometric}
\usepackage[hidelinks]{hyperref}
\setlength{\parindent}{0pt}
\setlength{\parskip}{3pt}
\setlist{nosep,leftmargin=*}

% ---------- Lösungsschalter ----------
\newif\ifloesung
\ifdefined\mitloesung\loesungtrue\fi
\newcommand{\lsgfarbe}{\color{red!75!black}}

% ---------- Kopf- und Fußzeile (einheitliche Form) ----------
\newcommand{\lehrkraft}{Hau}
\newcommand{\themenbereich}{3. Graphen}
\newcommand{\abnummer}{}
\pagestyle{fancy}\fancyhf{}
\renewcommand{\headrulewidth}{0.4pt}
\fancyhead[L]{\small Klasse 11\\Datum:}
\fancyhead[C]{\small Themenbereich:\\\themenbereich}
\fancyhead[R]{\small \lehrkraft\ifloesung\\{\lsgfarbe\bfseries Lösung}\fi}
\fancyfoot[C]{\scriptsize edu-mrh.de\,·\,Informatik 11\,·\,Blatt \abnummer\ifloesung\ (Lösung)\fi\,·\,Seite \thepage}
\newcommand{\abtitel}[2]{\renewcommand{\abnummer}{#1}\begin{center}\Large\bfseries #1\quad\underline{#2}\end{center}\vspace{-1mm}}

% ---------- Boxen ----------
\newenvironment{aufgabe}[2][]{\begin{tcolorbox}[enhanced,breakable,colback=white,colframe=black,boxrule=0.7pt,arc=3mm,left=2mm,right=2mm,top=1.5mm,bottom=1.5mm,before skip=2.5mm,after skip=2.5mm]\textbf{Aufgabe #2}\ifstrempty{#1}{}{ \textit{#1}}\par\vspace{0.6mm}}{\end{tcolorbox}}
\newtcolorbox{merke}[1][Merke]{enhanced,breakable,colback=green!5,colframe=green!40!black,boxrule=0.8pt,arc=2mm,left=2mm,right=2mm,top=1.5mm,bottom=1.5mm,before skip=2.5mm,after skip=2.5mm,title={#1},fonttitle=\bfseries,coltitle=white}
\newtcolorbox{hinweis}{enhanced,breakable,colback=yellow!12,colframe=orange!70!black,boxrule=0.6pt,arc=2mm,left=2mm,right=2mm,top=1mm,bottom=1mm,before skip=2mm,after skip=2mm}
\newtcolorbox{ausblick}[1][Ausblick]{enhanced,breakable,colback=teal!5,colframe=teal!60!black,boxrule=0.8pt,arc=2mm,left=2mm,right=2mm,top=1.5mm,bottom=1.5mm,before skip=2.5mm,after skip=2.5mm,title={#1},fonttitle=\bfseries,coltitle=white}
\newcommand{\web}[1]{{\small\color{gray}\textit{Website: #1}}}

% ---------- Lücken, Linien, Karos ----------
\newcommand{\luecke}[1][3.5cm]{\rule[-0.5ex]{#1}{0.4pt}}
\newcommand{\linien}[1]{\par\vspace{1.5mm}\foreach \i in {1,...,#1}{\noindent\rule[0pt]{\linewidth}{0.3pt}\par\vspace{5.5mm}}\vspace{-3mm}}
\newcommand{\kariert}[1]{\par\vspace{1.5mm}\noindent\begin{tikzpicture}\draw[step=5mm,gray!55,thin] (0,0) grid (\linewidth-1pt,#1*5mm);\end{tikzpicture}\par}
\newcommand{\karobox}[2]{\begin{tikzpicture}\draw[step=5mm,gray!55,thin] (0,0) grid (#1,#2);\end{tikzpicture}}
\newcommand{\code}[1]{\texttt{#1}}

% Lücke mit Lösung: \lsg[Mindestbreite]{Lösungstext}
\newlength{\lsgbreite}
\newcommand{\lsg}[2][3.5cm]{\ifloesung\settowidth{\lsgbreite}{\,#2\,}\ifdim\lsgbreite<#1\setlength{\lsgbreite}{#1}\fi\underline{\makebox[\lsgbreite][l]{\lsgfarbe\,#2}}\else\luecke[#1]\fi}
% Schreiblinien mit Lösung: \lsglinien{Anzahl}{Lösungstext} (gleiche Höhe wie \linien)
\newcommand{\lsglinien}[2]{\ifloesung\par\vspace{1.5mm}\noindent\begin{minipage}[t][#1\dimexpr 6.9mm\relax][t]{\linewidth}\lsgfarbe #2\end{minipage}\par\vspace{-1.5mm}\else\linien{#1}\fi}
% Beliebiger Inhalt: \lsgoder{nur in der Lösung}{nur im Blatt}
\newcommand{\lsgoder}[2]{\ifloesung\leavevmode#1\else#2\fi}

% ---------- Ergänzung (aus informatik/10/java/arbeitsblaetter/ab-vorlage.tex) ----------
% Bestehende Makros oben bleiben unverändert, damit ältere Blätter gleich aussehen.
%  \restzeile / \lsgrest{Text}   Lücke bis zum Zeilenende
%  \lsgtext{Text}                nur in der Lösung sichtbar (z. B. Tabellenzellen)
%  \lueckentext                  größerer Zeilenabstand, Flattersatz für Absätze mit Lücken
%  \schreibzeile                 Strut für Tabellenzellen zum Ausfüllen (8 mm hoch)
%  \alt{9}, \neu{7}              überschriebener alter Wert / neuer Wert (Dijkstra-Tabellen)
\usepackage{colortbl,xspace}
\usepackage[normalem]{ulem}
\newcommand{\lsgrest}[1]{\leavevmode\unskip\hspace{1.5mm}\ifloesung\rlap{\raisebox{0.8pt}{\lsgfarbe\,#1}}\fi\leaders\hbox{\rule[-0.6ex]{0.5pt}{0.4pt}}\hfill\null\par}
\newcommand{\restzeile}{\lsgrest{}}
\newcommand{\lsgtext}[1]{\ifloesung\leavevmode{\lsgfarbe#1}\fi}
\newcommand{\lueckentext}{\linespread{1.5}\selectfont\raggedright}
\newcommand{\schreibzeile}{\rule[-2.8mm]{0pt}{8mm}}
\newcommand{\alt}[1]{\leavevmode{\color{blue!60!black}\sout{#1}}}
\newcommand{\neu}[1]{\leavevmode{\color{green!50!black}\bfseries #1}}
\colorlet{djaktuell}{orange!40}
\colorlet{djfertig}{blue!22}

% ---------- Graphen zeichnen ----------
\tikzset{
  knoten/.style={draw,rounded corners=3pt,fill=orange!12,minimum height=6mm,inner sep=3pt,font=\small},
  kknoten/.style={draw,circle,fill=orange!12,minimum size=7mm,inner sep=1pt,font=\small},
  kante/.style={thick},
  gkante/.style={thick,-{Stealth[length=2.6mm]}},
  dkante/.style={thick,{Stealth[length=2.6mm]}-{Stealth[length=2.6mm]}},
  gewicht/.style={font=\scriptsize,fill=white,inner sep=1pt,text=teal!60!black},
  lsgkante/.style={thick,red!75!black},
  lsggkante/.style={thick,red!75!black,-{Stealth[length=2.6mm]}},
  lsgknoten/.style={knoten,draw=red!75!black,text=red!75!black,fill=red!4},
}
% Adjazenzmatrix zum Ausfüllen: \adjmx[Lösungseinträge]{Anzahl}{Knotenliste}{Zellbreite}
% Lösungseinträge als Zeile/Spalte/Text, z. B. \adjmx[1/2/X,2/1/X]{3}{A,B,C}{7mm}
\newcommand{\adjmx}[4][]{\begin{tikzpicture}[x=#4,y=#4]
  \foreach \n [count=\i] in {#3} {\node[font=\small\bfseries] at (\i,0) {\n}; \node[font=\small\bfseries,anchor=east] at (0.45,-\i) {\n};}
  \foreach \k in {0,...,#2} {\draw[gray] (\k+0.5,-0.5) -- (\k+0.5,-#2-0.5); \draw[gray] (0.5,-\k-0.5) -- (#2+0.5,-\k-0.5);}
  \ifloesung\ifstrempty{#1}{}{\foreach \z/\s/\t in {#1} {\node[text=red!75!black,font=\small] at (\s,-\z) {\t};}}\fi
\end{tikzpicture}}

% ---------- Java-Listings ----------
\lstset{language=Java,basicstyle=\ttfamily\small,keywordstyle=\bfseries,commentstyle=\color{gray!80!black},stringstyle=\color{black},showstringspaces=false,columns=flexible,keepspaces=true,tabsize=3,frame=single,framerule=0.4pt,rulecolor=\color{black},xleftmargin=2mm,framexleftmargin=1mm,aboveskip=2mm,belowskip=2mm,breaklines=true,escapechar=§,
 literate={ä}{{\"a}}1 {ö}{{\"o}}1 {ü}{{\"u}}1 {Ä}{{\"A}}1 {Ö}{{\"O}}1 {Ü}{{\"U}}1 {ß}{{\ss}}1 {–}{{--}}1 {„}{{\glqq}}1 {“}{{\grqq}}1}
\lstnewenvironment{java}[1][]{\lstset{#1}}{}

% Dörfer-Graph (Merkkasten Begriffe): Koordinaten wie im Begriffe-Explorer der Website
\newcommand{\doerferknoten}{%
  \node[knoten,font=\tiny,minimum height=4mm,inner sep=1.5pt] (E) at (1.2,3.15) {Eulenhain};
  \node[knoten,font=\tiny,minimum height=4mm,inner sep=1.5pt] (S) at (3.3,3.25) {Sternwies};
  \node[knoten,font=\tiny,minimum height=4mm,inner sep=1.5pt] (K) at (5.2,2.65) {Kieselfurt};
  \node[knoten,font=\tiny,minimum height=4mm,inner sep=1.5pt] (F) at (2.5,1.85) {Funkenau};
  \node[knoten,font=\tiny,minimum height=4mm,inner sep=1.5pt] (H) at (4.8,1.25) {Haselbruck};
  \node[knoten,font=\tiny,minimum height=4mm,inner sep=1.5pt] (N) at (0.9,0.75) {Nebelried};
  \node[knoten,font=\tiny,minimum height=4mm,inner sep=1.5pt] (T) at (3.2,0.35) {Tannzapf};}
\begin{document}
\abtitel{3.1}{Graphen -- Grundbegriffe}

Hänsel und Gretel haben die Hexe überlistet. Doch die Brotkrumen, mit denen sie ihren Weg markiert hatten, haben die Vögel gefressen. Findest du trotzdem nach Hause?

\begin{aufgabe}[Hänsel und Gretel]{1}
Starte das Spiel auf der Website (Station 3.1).
\textbf{a)} Zeichne beim Spielen eine Karte des Waldes: jeder Ort ein Kästchen, jeder Weg eine Linie. Führt ein Weg nur in eine Richtung, zeichne einen Pfeil. Das Hexenhaus ist schon eingezeichnet.
\end{aufgabe}
\noindent\begin{tikzpicture}
\draw[step=5mm,gray!55,thin] (0,0) grid (17.4,7);
\node[knoten,fill=green!15] (hex) at (4,2.3) {Hexenhaus (Start)};
\ifloesung
  \node[lsgknoten] (warn) at (5.5,6.5) {Warnung};
  \node[lsgknoten] (wolf) at (10.5,6.5) {Werwolf};
  \node[lsgknoten] (stadt) at (1.6,5.2) {Stadt};
  \node[lsgknoten] (abz) at (5.5,5.2) {Abzweigung};
  \node[lsgknoten] (spinne) at (13,5.2) {Spinne};
  \node[lsgknoten] (baeck) at (1.6,3.6) {Bäckerei};
  \node[lsgknoten] (pilze) at (4,3.9) {Pilze};
  \node[lsgknoten] (w3) at (13,3.4) {Wald3};
  \node[lsgknoten] (w1) at (8,1.0) {Wald1};
  \node[lsgknoten] (stein) at (12.2,2.0) {Steinkreis};
  \node[lsgknoten] (w2) at (15,0.8) {Wald2};
  \node[lsgknoten] (heim) at (16.2,2.6) {Zuhause};
  \draw[lsgkante] (stadt) -- (baeck);
  \draw[lsgkante] (stadt) -- (pilze);
  \draw[lsgkante] (pilze) -- (abz);
  \draw[lsgkante] (pilze) -- (hex);
  \draw[lsggkante] (pilze) to[bend left=50] node[right,font=\scriptsize,text=red!75!black] {Pilze essen} (hex);
  \draw[lsgkante] (abz) -- (warn);
  \draw[lsggkante] (warn) -- (wolf);
  \draw[lsgkante] (abz) -- (spinne);
  \draw[lsgkante] (spinne) -- (w3);
  \draw[lsgkante] (hex) -- (w1);
  \draw[lsgkante] (w1) -- (w3);
  \draw[lsgkante] (w1) -- (w2);
  \draw[lsgkante] (w1) -- (stein);
  \draw[lsgkante] (w2) -- (w3);
  \draw[lsgkante] (w2) -- (stein);
  \draw[lsgkante] (w3) -- (stein);
  \draw[lsggkante] (w2) -- (heim);
  \node[text=red!75!black,font=\scriptsize,align=left,anchor=north west] at (0.2,1.3) {13 Knoten, 17 Kanten\\(14 ungerichtet, 3 gerichtet)};
\fi
\end{tikzpicture}

\textbf{b)} Am Ende zeigt dir das Spiel deinen Weg. Notiere ihn und die Anzahl der Schritte:
\lsglinien{1}{individuell; kürzester Weg: Hexenhaus -- Wald1 -- Wald2 -- Zuhause (3 Schritte)}

\begin{minipage}[t]{0.6\linewidth}
\begin{merke}[Graph und Kurzschreibweise $G = (V, E)$]\raggedright
Ein Graph besteht aus \lsg[2.3cm]{Knoten} (engl. \textit{vertices}) und \lsg[2.3cm]{Kanten} (engl. \textit{edges}). Eine Kante verbindet zwei Knoten. Im Spiel sind die Orte die \lsg[2cm]{Knoten}, die Wege die \lsg[2cm]{Kanten}.\\[1mm]
Kurzschreibweise $G = (V, E)$: $V$ ist die Menge der Knoten, $E$ die Menge der Kanten.
\begin{itemize}
\item $(A, B)$: \lsg[2.5cm]{gerichtete} Kante, nur von $A$ nach $B$
\item $\{A, B\}$: \lsg[2.5cm]{ungerichtete} Kante, steht für $(A, B)$ und \lsg[1.4cm]{$(B, A)$}
\end{itemize}
\end{merke}
\end{minipage}\hfill
\begin{minipage}[t]{0.37\linewidth}
\vspace{1mm}
\centering\small \textbf{Beispiel: Wer folgt wem?}\\[1mm]
\begin{tikzpicture}[x=0.95cm,y=0.9cm]
\node[knoten] (mira) at (2.4,2.5) {Mira};
\node[knoten] (ole) at (4.3,1.9) {Ole};
\node[knoten] (zeynep) at (3.4,0.3) {Zeynep};
\node[knoten] (finn) at (0.9,1.0) {Finn};
\node[knoten] (paula) at (0.5,2.5) {Paula};
\draw[dkante] (mira) -- (ole);
\draw[gkante] (ole) -- (zeynep);
\draw[gkante] (zeynep) -- (mira);
\draw[dkante] (mira) -- (finn);
\draw[gkante] (finn) -- (zeynep);
\end{tikzpicture}
\end{minipage}\\[1mm]
{\small $V = \{\text{Finn}, \text{Mira}, \text{Ole}, \text{Paula}, \text{Zeynep}\}$\\
$E = \{\{\text{Mira}, \text{Ole}\}, \{\text{Mira}, \text{Finn}\}, (\text{Ole}, \text{Zeynep}), (\text{Zeynep}, \text{Mira}), (\text{Finn}, \text{Zeynep})\}$}

\begin{aufgabe}[Hänsel und Gretel]{1}
\textbf{c)} Schreibe den Ausschnitt deiner Karte mit den Knoten Hexenhaus, Pilze, Stadt, Bäckerei und Wald1 in Kurzschreibweise auf (5 Kanten).\\[2mm]
$V = \{$\lsg[15.4cm]{Hexenhaus, Pilze, Stadt, Bäckerei, Wald1}$\}$\\[3mm]
$E = \{$\lsg[15.4cm]{\{Hexenhaus, Pilze\}, (Pilze, Hexenhaus), \{Pilze, Stadt\}, \{Stadt, Bäckerei\}, \{Hexenhaus, Wald1\}}$\}$
\end{aufgabe}

\newpage
\begin{minipage}[t]{0.71\linewidth}
\begin{merke}[Pfade, Zyklen, Richtung, Gewicht]\raggedright\small
\textbf{Pfad (Weg):} Folge von Knoten, in der aufeinanderfolgende Knoten durch eine \lsg[1.8cm]{Kante} verbunden sind; keine Kante wird doppelt benutzt.\\[0.6mm]
\textbf{Einfacher Pfad:} Pfad, in dem sich kein \lsg[1.8cm]{Knoten} wiederholt.\\
Beispiel in A: \lsg[5.6cm]{Nebelried -- Funkenau -- Kieselfurt -- Haselbruck}\\[0.6mm]
\textbf{Länge eines Pfades:} ungewichtet die \lsg[3.2cm]{Anzahl der Kanten}, gewichtet die \lsg[3.2cm]{Summe der Gewichte}.\\[0.6mm]
\textbf{Zyklus:} einfacher Pfad, dessen \lsg[5cm]{Start- und Endknoten gleich sind}.\\
Beispiel in A: \lsg[5.6cm]{Sternwies -- Funkenau -- Kieselfurt -- Sternwies}\\
Graph mit Zyklus: \emph{zyklisch}, sonst \lsg[3.4cm]{azyklisch (zyklenfrei)}.\\[0.6mm]
\textbf{Gerichtet:} Kanten nur \lsg[3cm]{in Pfeilrichtung} (Pfeil).\ \textbf{Ungerichtet:} \lsg[2.8cm]{in beide Richtungen} (Strich).\\[0.6mm]
\textbf{Stark zusammenhängend:} Von \lsg[2.4cm]{jedem Knoten} gibt es einen Pfad zu \lsg[3.2cm]{jedem anderen Knoten}.\\
\textbf{Schwach zusammenhängend:} \lsg[5.4cm]{nur ohne Beachtung der Richtungen}\\ \lsg[\linewidth]{hängt alles zusammen. B ist schwach: Von Tannzapf geht keine Kante weg.}\\[0.6mm]
\textbf{Gewichtet:} Jede Kante trägt eine \lsg[2.6cm]{Zahl (Gewicht)}, z.\,B. \lsg[3.6cm]{Entfernung, Fahrzeit}.
\end{merke}
\end{minipage}\hfill
\begin{minipage}[t]{0.27\linewidth}
\vspace{1mm}
\textbf{\small A} {\scriptsize ungerichtet}\\
\begin{tikzpicture}[x=0.8cm,y=0.85cm]
\doerferknoten
\draw[kante] (E)--(S); \draw[kante] (E)--(F); \draw[kante] (S)--(F); \draw[kante] (S)--(K);
\draw[kante] (F)--(K); \draw[kante] (F)--(N); \draw[kante] (K)--(H); \draw[kante] (H)--(F);
\draw[kante] (N)--(T); \draw[kante] (H)--(T);
\end{tikzpicture}\\[2mm]
\textbf{\small B} {\scriptsize gerichtet}\\
\begin{tikzpicture}[x=0.8cm,y=0.85cm]
\doerferknoten
\draw[gkante] (E)--(S); \draw[gkante] (E)--(F); \draw[gkante] (S)--(F); \draw[gkante] (S)--(K);
\draw[gkante] (F)--(K); \draw[gkante] (F)--(N); \draw[gkante] (K)--(H); \draw[gkante] (H)--(F);
\draw[gkante] (N)--(T); \draw[gkante] (H)--(T);
\end{tikzpicture}
\end{minipage}

\begin{aufgabe}{2}
{\raggedright Zeichne $G = (V, E)$ mit $V = \{A, B, C, D, E, F\}$ und\\ $E = \{(A, B), (A, C), (A, E), (B, C), (C, C), (C, D), (E, C), (F, D)\}$. Beantworte die Fragen.\par}\vspace{1mm}
\begin{minipage}[t]{0.42\linewidth}
\vspace{0pt}
\begin{tikzpicture}
\draw[step=5mm,gray!55,thin] (0,0) grid (7,3);
\ifloesung
  \foreach \n/\x/\y/\t in {a/0.6/1.5/A, b/2/2.55/B, e/2/0.45/E, c/3.4/1.5/C, d/4.9/1.5/D, f/6.4/1.5/F} {
    \node[kknoten,minimum size=6mm,draw=red!75!black,text=red!75!black] (\n) at (\x,\y) {\t};}
  \draw[lsggkante] (a)--(b); \draw[lsggkante] (a)--(c); \draw[lsggkante] (a)--(e);
  \draw[lsggkante] (b)--(c); \draw[lsggkante] (c)--(d); \draw[lsggkante] (e)--(c);
  \draw[lsggkante] (f)--(d);
  \draw[lsggkante] (c) to[out=60,in=120,looseness=6] (c);
\fi
\end{tikzpicture}
\end{minipage}\hfill
\begin{minipage}[t]{0.55\linewidth}
\vspace{0pt}\small\raggedright
Gerichtet? Woran erkennbar? \lsg[3.9cm]{ja, Paare ( ) statt \{ \}}\\[1.5mm]
Ein Zyklus: \lsg[5.7cm]{(C, C), eine Schlinge (Länge 1)}\\[1.5mm]
Stark/schwach zusammenhängend? \lsg[3.6cm]{schwach: von D}\\[1.5mm]
\lsg[0.98\linewidth]{kommt man nirgends hin.}\\[1.5mm]
Länge des längsten einfachen Pfades: \lsg[3.2cm]{3 (A--B--C--D)}
\end{minipage}
\end{aufgabe}

\begin{aufgabe}[Zurück in den Wald]{3}
Deine Karte aus Aufgabe 1: \textbf{a)} gerichtet?\ \textbf{b)} ein Zyklus?\ \textbf{c)} stark oder schwach zusammenhängend?\ \textbf{d)} „Der Weg zurück … scheint verschwunden“: Ist die Kante Hexenhaus -- Wald1 gerichtet? Prüfe im Spiel.
\lsglinien{2}{\small a) gemischt; gerichtet nur (Pilze, Hexenhaus), (Warnung, Werwolf), (Wald2, Zuhause)\quad b) Wald1 -- Wald2 -- Wald3 -- Wald1\\ c) schwach: Von Zuhause/Werwolf führt keine Kante weg.\quad d) nein: Von Wald2 aus führt ein Weg zurück.}
\end{aufgabe}

\begin{aufgabe}[Ordne die Graphen ein.]{4}
\vspace{-1mm}
\small\renewcommand{\arraystretch}{1.15}
\begin{tabularx}{\linewidth}{|>{\centering\arraybackslash}m{5.6cm}|>{\centering\arraybackslash}m{1.7cm}|>{\centering\arraybackslash}m{1.1cm}|>{\centering\arraybackslash}m{1.1cm}|>{\centering\arraybackslash}m{1.1cm}|>{\centering\arraybackslash}X|}\hline
\textbf{Graph} & \scriptsize\textbf{zusammen-\newline hängend?}\newline stark/schwach/nein & \scriptsize\textbf{ge-\newline richtet?} & \scriptsize\textbf{ge-\newline wichtet?} & \scriptsize\textbf{zyklisch?} & \scriptsize\textbf{Länge eines längsten einfachen Pfades}\\\hline
\rule{0pt}{8mm}\textbf{I}\ \begin{tikzpicture}[baseline=(k5.base),x=0.85cm,y=0.55cm]
\foreach \n/\x/\y in {1/0/1,2/1/1,3/2/1,4/3/1,5/1/0,6/3/0} {\node[kknoten,minimum size=5mm,inner sep=0pt,font=\tiny] (k\n) at (\x,\y) {K\n};}
\draw[kante] (k1)--(k2)--(k3)--(k4); \draw[kante] (k2)--(k5);
\end{tikzpicture}
& \lsgoder{\lsgfarbe nein}{} & \lsgoder{\lsgfarbe nein}{} & \lsgoder{\lsgfarbe nein}{} & \lsgoder{\lsgfarbe nein}{} & \lsgoder{\lsgfarbe 3: K1--K2--K3--K4}{}\\\hline
\rule{0pt}{5mm}\textbf{II}\ $V = \{A, B, C\}$\newline $E =\{(A, B), (B, A), (C, B), (C, C)\}$
& \lsgoder{\lsgfarbe schwach}{} & \lsgoder{\lsgfarbe ja}{} & \lsgoder{\lsgfarbe nein}{} & \lsgoder{\lsgfarbe ja}{} & \lsgoder{\lsgfarbe 2: C$\to$B$\to$A}{}\\\hline
\rule{0pt}{8mm}\textbf{III}\ \begin{tikzpicture}[baseline=(k3.base),x=0.9cm,y=0.6cm]
\foreach \n/\x/\y in {1/0/1,2/2/1,3/1/0,4/3/0} {\node[kknoten,minimum size=5mm,inner sep=0pt,font=\tiny] (k\n) at (\x,\y) {K\n};}
\draw[kante] (k1) -- node[gewicht] {4} (k2); \draw[kante] (k2) -- node[gewicht] {1} (k3);
\draw[kante] (k3) -- node[gewicht] {5} (k1); \draw[kante] (k3) -- node[gewicht] {2} (k4);
\end{tikzpicture}
& \lsgoder{\lsgfarbe stark}{} & \lsgoder{\lsgfarbe nein}{} & \lsgoder{\lsgfarbe ja}{} & \lsgoder{\lsgfarbe ja}{} & \lsgoder{\lsgfarbe 11: K2--K1--K3--K4}{}\\\hline
\end{tabularx}
\end{aufgabe}
\end{document}
